% GENERATED FILE. Edit canonical lessons, then run npm run generate:books.
% canonical-source-hash: fnv1a64:9bbc01112d391574
% canonical-lessons: ES-C777-congelar, ES-C777-enfriar, ES-C777-planchar, ES-C777-enchufar, ES-C777-reiniciar, ES-R777-repaso-776-777, ES-R777-repaso-final-777

\chapter{Si se va la luz}
\label{ch:es-si-se-va-la-luz}
\hlchaptermodality{777}

\begin{chapteropening}
\textbf{By the end of this chapter:} \emph{I can say what to do if the power goes or food goes off: freeze it, let it cool, iron later, plug in and restart.}

Complete the last lesson of chapter 777: I can say what to do if the power goes or food goes off: freeze it, let it cool, iron later, plug in and restart.
\end{chapteropening}

\section[congelar]{congelar — to freeze}

\label{lesson:ES-C777-congelar}



\begin{quote}
Before the new one: say the Spanish for to update, then the Spanish for the keyboard.
\end{quote}

\subsection*{You'll want to know: congelar}
\textbf{congelar} — \textquotedblleft{}to freeze\textquotedblright{}.

\begin{quote}
Si no vas a comerlo hoy, congélalo.
\end{quote}

\begin{quote}
\emph{If you aren't going to eat it today, freeze it.}
\end{quote}

\emph{Congelado} — frozen. \emph{Los congelados} — frozen food.

\begin{cousinweb}
\emph{congelar} is a learned borrowing from Latin \emph{congelāre} 'to cause to freeze' (EtymDB), built on \emph{gelū} 'frost', the word that gave \emph{hielo}.
\end{cousinweb}

\begin{grammarlens}[title={What you've built}]
One more word for telling, explaining and describing at B1.
\end{grammarlens}

\subsection*{Your turn}
\begin{itemize}
\raggedright
  \item \emph{Say it:} \emph{congelar}
  \item \emph{Say it:} \emph{congelar}, once more
  \item \emph{Say it:} a sentence of your own with \emph{congelar}
  \item \emph{Recall:} say the Spanish for to update, then the Spanish for the keyboard, then say \emph{congelar} again
\end{itemize}

\begin{tcolorbox}[breakable,colback=teal!4,colframe=teal!35!black,title={Before you move on}]
What does congelar mean? (\textquotedblleft{}To freeze\textquotedblright{}.) Say it once more.
\end{tcolorbox}
\section[enfriar]{enfriar — to cool, to chill}

\label{lesson:ES-C777-enfriar}



\begin{quote}
Before the new one: say the Spanish for the keyboard, then the Spanish for to freeze.
\end{quote}

\subsection*{You'll want to know: enfriar}
\textbf{enfriar} — \textquotedblleft{}to cool, to chill\textquotedblright{}.

\begin{quote}
Si la sopa está muy caliente, déjala enfriar.
\end{quote}

\begin{quote}
\emph{If the soup is very hot, let it cool.}
\end{quote}

It takes an accent: \emph{enfrío}. \emph{Enfriarse} is to cool down, and also to catch a cold.

\begin{cousinweb}
\emph{enfriar} is built on \emph{frío} 'cold' with \emph{en-} and \emph{-ar}. \emph{Frío} is inherited from Latin \emph{frīgidus} 'cold' (Lewis \& Short).
\end{cousinweb}

\begin{grammarlens}[title={What you've built}]
One more word for telling, explaining and describing at B1.
\end{grammarlens}

\subsection*{Your turn}
\begin{itemize}
\raggedright
  \item \emph{Say it:} \emph{enfriar}
  \item \emph{Say it:} \emph{enfriar}, once more
  \item \emph{Say it:} a sentence of your own with \emph{enfriar}
  \item \emph{Recall:} say the Spanish for the keyboard, then the Spanish for to freeze, then say \emph{enfriar} again
\end{itemize}

\begin{tcolorbox}[breakable,colback=teal!4,colframe=teal!35!black,title={Before you move on}]
What does enfriar mean? (\textquotedblleft{}To cool, to chill\textquotedblright{}.) Say it once more.
\end{tcolorbox}
\section[planchar]{planchar — to iron}

\label{lesson:ES-C777-planchar}



\begin{quote}
Before the new one: say the Spanish for to freeze, then the Spanish for to cool.
\end{quote}

\subsection*{You'll want to know: planchar}
\textbf{planchar} — \textquotedblleft{}to iron\textquotedblright{}.

\begin{quote}
Si tienes tiempo, plancha la camisa.
\end{quote}

\begin{quote}
\emph{If you have time, iron the shirt.}
\end{quote}

\emph{La plancha} is the iron; \emph{a la plancha} is grilled on a hot plate.

\begin{cousinweb}
\emph{planchar} is built on \emph{plancha} 'a flat iron plate', which comes from French \emph{planche} 'a board' (EtymDB), from Latin \emph{planca} 'a plank' (Lewis \& Short).
\end{cousinweb}

\begin{grammarlens}[title={What you've built}]
One more word for telling, explaining and describing at B1.
\end{grammarlens}

\subsection*{Your turn}
\begin{itemize}
\raggedright
  \item \emph{Say it:} \emph{planchar}
  \item \emph{Say it:} \emph{planchar}, once more
  \item \emph{Say it:} a sentence of your own with \emph{planchar}
  \item \emph{Recall:} say the Spanish for to freeze, then the Spanish for to cool, then say \emph{planchar} again
\end{itemize}

\begin{tcolorbox}[breakable,colback=teal!4,colframe=teal!35!black,title={Before you move on}]
What does planchar mean? (\textquotedblleft{}To iron\textquotedblright{}.) Say it once more.
\end{tcolorbox}
\section[enchufar]{enchufar — to plug in}

\label{lesson:ES-C777-enchufar}



\begin{quote}
Before the new one: say the Spanish for to cool, then the Spanish for to iron.
\end{quote}

\subsection*{You'll want to know: enchufar}
\textbf{enchufar} — \textquotedblleft{}to plug in\textquotedblright{}.

\begin{quote}
Si no carga, enchufa el cable otra vez.
\end{quote}

\begin{quote}
\emph{If it isn't charging, plug the cable in again.}
\end{quote}

\emph{El enchufe} is the plug or socket; in Spain it also means connections, knowing the right people.

\begin{cousinweb}
The origin of \emph{enchufar} is not settled. It is usually treated as imitative, built on a sound like \emph{chuf}, the noise of one pipe pushed into another.
\end{cousinweb}

\begin{grammarlens}[title={What you've built}]
One more word for telling, explaining and describing at B1.
\end{grammarlens}

\subsection*{Your turn}
\begin{itemize}
\raggedright
  \item \emph{Say it:} \emph{enchufar}
  \item \emph{Say it:} \emph{enchufar}, once more
  \item \emph{Say it:} a sentence of your own with \emph{enchufar}
  \item \emph{Recall:} say the Spanish for to cool, then the Spanish for to iron, then say \emph{enchufar} again
\end{itemize}

\begin{tcolorbox}[breakable,colback=teal!4,colframe=teal!35!black,title={Before you move on}]
What does enchufar mean? (\textquotedblleft{}To plug in\textquotedblright{}.) Say it once more.
\end{tcolorbox}
\section[reiniciar]{reiniciar — to restart, to reboot}

\label{lesson:ES-C777-reiniciar}



\begin{quote}
Before the new one: say the Spanish for to iron, then the Spanish for to plug in.
\end{quote}

\subsection*{You'll want to know: reiniciar}
\textbf{reiniciar} — \textquotedblleft{}to restart, to reboot\textquotedblright{}.

\begin{quote}
Si se bloquea, reinicia el móvil.
\end{quote}

\begin{quote}
\emph{If it freezes, restart your phone.}
\end{quote}

\begin{cousinweb}
\emph{reiniciar} is built on \emph{iniciar} with \emph{re-} 'again'. \emph{Iniciar} is a learned borrowing from Latin \emph{initiāre} 'to begin', from \emph{initium} 'a beginning' (Lewis \& Short).
\end{cousinweb}

\begin{grammarlens}[title={What you've built}]
One more word for telling, explaining and describing at B1.
\end{grammarlens}

\subsection*{Your turn}
\begin{itemize}
\raggedright
  \item \emph{Say it:} \emph{reiniciar}
  \item \emph{Say it:} \emph{reiniciar}, once more
  \item \emph{Say it:} a sentence of your own with \emph{reiniciar}
  \item \emph{Recall:} say the Spanish for to iron, then the Spanish for to plug in, then say \emph{reiniciar} again
\end{itemize}

\begin{tcolorbox}[breakable,colback=teal!4,colframe=teal!35!black,title={Before you move on}]
What does reiniciar mean? (\textquotedblleft{}To restart, to reboot\textquotedblright{}.) Say it once more.
\end{tcolorbox}
\section[(dialogue)]{Más cosas de cada día}

\label{lesson:ES-R777-repaso-776-777}



\begin{quote}
Say the Spanish for the password and the file.
\end{quote}

\subsection*{Your turn}
Say these aloud:

\begin{itemize}
\raggedright
  \item the password, the file, the folder, to update, the keyboard
  \item to freeze, to cool, to iron, to plug in, to restart
\end{itemize}

\begin{tcolorbox}[breakable,colback=teal!4,colframe=teal!35!black,title={Before you move on}]
The password? (\textbf{la contraseña}.) The file? (\textbf{el archivo}.) The folder? (\textbf{la carpeta}.) To update? (\textbf{actualizar}.) The keyboard? (\textbf{el teclado}.) To freeze? (\textbf{congelar}.) To cool? (\textbf{enfriar}.) To iron? (\textbf{planchar}.) To plug in? (\textbf{enchufar}.) To restart? (\textbf{reiniciar}.)
\end{tcolorbox}
\section[(dialogue)]{Si se va la luz — once more}

\label{lesson:ES-R777-repaso-final-777}



\begin{quote}
Say the Spanish for to freeze and to cool.
\end{quote}

\subsection*{Your turn}
\begin{itemize}
\raggedright
  \item \emph{Say it:} to freeze, to cool, to iron, to plug in, to restart
\end{itemize}

\begin{tcolorbox}[breakable,colback=teal!4,colframe=teal!35!black,title={Before you move on}]
To freeze? (\textbf{congelar}.) To cool? (\textbf{enfriar}.) To iron? (\textbf{planchar}.) To plug in? (\textbf{enchufar}.) To restart? (\textbf{reiniciar}.)
\end{tcolorbox}
